% ============================================================
\section{The Bedrock Principles}
\label{sec:diophantine}
% ============================================================

The monograph begins from one primitive, three structural axioms, and two
dynamic bedrock laws. The primitive, Constancy, asserts only that things can
be told apart through decidable identity. Existence, Incidence, and Closure
Balance force a finite compositional counting grammar. The termination and
exactification laws added below then identify exact closure and state what a
non-closed incidence must do. Crucially, no
metric, topology, continuum parameter, constitutive law, or physical time is
imposed in advance. This absence of imposed structure---not any
presence---is what forces the combinatorial closure grammar and exactification
spine developed below.

This manuscript is not presented as a discrete spacetime ontology or as a
replacement for other fundamental-physics programs. Its narrower question is:
under finite observability, when does a restricted support carry an autonomous
law of its own? The combinatorics of Parts~I--II provide the minimal
exactification grammar for that question; geometry, continuum time, and the
flagship equations enter only later as selected autonomous realizations or
representative normal forms under explicit additional hypotheses.

The ``no further structure'' clause deserves emphasis. Beyond the explicit
incidence, closure, termination, and exactification interfaces introduced in
this chapter, no metric, topology, algebraic operation, continuum parameter,
or external selection rule is given. On that surface the manuscript adopts the
\emph{complete} simplicial complex on the available primitives as the
canonical no-further-structure closure convention: if a face were excluded,
that exclusion would amount to additional structural datum beyond the present
axiom surface. Within that convention, together with the finite combinatorics
of the next chapter and the later soundness bridge, the closure equation
selects the canonical nondegenerate balanced carrier~$D=4$ and the carrier
decomposition~$1{+}2{+}3$. The
concept of decidable equality in a constructive setting is standard
\cite{Bishop1967,MartinLof1984,deMouraKongAvigadEtAl2015}; what is new here is
the derivation chain that begins from it.

% ------------------------------------------------------------
\subsection{Primitive and Structural Axioms}
\label{subsec:three-axioms}
% ------------------------------------------------------------

\begin{primitive}[Constancy]\lean{Primitive}
A \emph{constant object} is a bearer of decidable self-identity: given any two
instances, their equality is decidable. Constancy records that the object
persists as the same self-contained whole until incidence with another object
forces a new composed description. Since time is not yet defined at this
stage, Constancy is not ``unchanging in time'' but decidable identity prior to
any temporal interpretation.
\end{primitive}

In particular, a function may be constant in this primitive sense without being
pointwise constant: the constancy lies in the persistence of the same identity
as a whole.

\begin{definition}[Primitive]\label{def:primitive}\lean{Primitive}
A \emph{primitive object}, or simply a \emph{primitive}, is an irreducible bearer of Constancy.
\end{definition}

\begin{axiom}[Existence]\label{ax:extent}\lean{Primitive}
There exists a primitive.
Primitives are taken with a determinate identity relation, so distinction can
be formulated at the counting level itself.
\end{axiom}

\begin{definition}[Incidence]\label{def:incidence}\lean{Event}
An \emph{incidence} is a finite composite formed by incident primitives.
\end{definition}

\begin{axiom}[Incidence]\label{ax:meeting}\lean{Event}
Primitives admit incidence.
At this stage an incidence carries at least its finite dimension datum and the
subincidence/presentation structure needed for the exactification law below.
The counting surface of the next chapter uses only the dimension datum.
\end{axiom}

\begin{axiom}[Closure Balance]\label{ax:closure}\lean{bulkCount}\lean{boundaryCount}\lean{closure}
Every complete count must close: the number of admissible pair counts of a structure equals the number of admissible order counts of that structure.
\end{axiom}

These three structural axioms generate arithmetic and a compositional counting
grammar. The termination and exactification laws added below supply the
zero-defect endpoint and the continuation clause for nonterminal incidence.
Any physical interpretation belongs to a later realization theorem, not to the
bedrock itself.

% ------------------------------------------------------------
\subsection{Termination and Exactification}
\label{subsec:exactification-terminal-incidence}
% ------------------------------------------------------------

The counting axioms say when an incidence is closed, but by themselves they do
not yet say how that closure is recognized at bedrock or what a non-closed
incidence must do. The missing dynamic content is therefore twofold:
termination identifies exact closure by vanishing defect, and exactification
forces any nonterminal incidence to continue by strict extension. These laws
still do not introduce time. They state only how an incidence stands with
respect to exact closure and further composition.

\begin{definition}[Extension of incidences]
\label{def:incidence-extension}\lean{CoherenceCalculus.incidenceExtension}
Let \(I\) and \(J\) be incidences. Write
\[
I \preceq J
\]
and say that \(J\) is an \emph{extension} of \(I\) if \(J\) contains \(I\) as
a subincidence and preserves every already-incident primitive relation of
\(I\). It is a \emph{strict extension}, written
\[
I \prec J,
\]
if \(I \preceq J\) and \(J\neq I\).
\end{definition}

\begin{proposition}[Extension is a preorder]
\label{prop:extension-is-preorder}\lean{CoherenceCalculus.extension_is_preorder}
The relation \(\preceq\) of
Definition~\ref{def:incidence-extension} is reflexive and transitive on
incidences.
\end{proposition}

\begin{proof}
Reflexivity is immediate since every incidence contains itself as a
subincidence and preserves all already-incident primitive relations of itself.
Transitivity follows because if \(I\preceq J\) and \(J\preceq K\), then
\(K\) contains \(I\) as a subincidence and preserves every primitive relation
already incident in \(I\).
\end{proof}

\begin{definition}[Lift of Constancy to an incidence]
\label{def:constancy-lift}\lean{CoherenceCalculus.ConstancyLift}
Let \(I\) be an incidence. A \emph{lift of Constancy to \(I\)} is a compatible
identity structure on the whole incidence \(I\) that:
\begin{enumerate}[label=(\roman*), leftmargin=2em, itemsep=0.2em]
\item agrees with the primitive identities already carried by the constituents
      of \(I\);
\item is stable under incidence-isomorphic presentations of the same
      incidence; and
\item preserves the identity of every subincidence already contained in \(I\).
\end{enumerate}
\end{definition}

\begin{definition}[Terminal incidence]
\label{def:terminal-incidence}\lean{CoherenceCalculus.TerminalIncidence}
An incidence \(I\) is \emph{terminal} if Constancy lifts to \(I\) as a whole in
the sense of Definition~\ref{def:constancy-lift}.
\end{definition}

\begin{proposition}[Terminal subincidence is inherited]
\label{prop:terminal-subincidence-inherited}\lean{CoherenceCalculus.terminal_subincidence_inherited}
If \(K\preceq I\preceq J\) and \(K\) is terminal, then the lifted identity
structure on \(K\) remains preserved in \(J\).
\end{proposition}

\begin{proof}
By Definition~\ref{def:incidence-extension}, extension preserves every
already-incident primitive relation of the smaller incidence. Since
terminality is the existence of a compatible Constancy lift on \(K\), that
lift remains preserved under further extension.
\end{proof}

\begin{definition}[Closure-defect space]
\label{def:closure-defect-space}\lean{CoherenceCalculus.ClosureDefectSpace}
A \emph{closure-defect space} is a pointed class \((\mathfrak D,0)\).
\end{definition}

\begin{definition}[Closure defect]
\label{def:closure-defect}\lean{CoherenceCalculus.ClosureDefectData, CoherenceCalculus.closureDefect}
Fix a closure-defect space \((\mathfrak D,0)\). The \emph{closure defect} of
an incidence \(I\) is a defect datum
\[
\Lambda(I)\in\mathfrak D.
\]
The distinguished point \(0\in\mathfrak D\) is the zero-defect class.
\end{definition}

\begin{axiom}[Termination]
\label{ax:termination}\lean{CoherenceCalculus.TerminationData, CoherenceCalculus.terminal_iff_zero_defect}
Fix a closure-defect space \((\mathfrak D,0)\) and a closure defect assignment
\(I\mapsto \Lambda(I)\in\mathfrak D\). Then for every incidence \(I\),
\[
I\text{ is terminal}
\quad\Longleftrightarrow\quad
\Lambda(I)=0.
\]
\end{axiom}

\begin{definition}[Exactifying continuation class]
\label{def:exactifying-continuation-class}\lean{CoherenceCalculus.ExactifyingContinuationClass}
For an incidence \(I\), define its \emph{exactifying continuation class} by
\[
\mathrm{Ext}(I)
:=
\{\,J\succ I : J \text{ preserves every already-incident primitive relation of } I\,\}.
\]
\end{definition}

\begin{axiom}[Nonretroactive exactification]
\label{ax:nonretroactive-exactification}\lean{CoherenceCalculus.NonretroactiveExactificationData}
Fix a closure-defect space \((\mathfrak D,0)\) and a closure defect assignment
\(I\mapsto \Lambda(I)\in\mathfrak D\). Then for every incidence \(I\):
\begin{enumerate}[label=(\roman*), leftmargin=2em, itemsep=0.2em]
\item if \(\Lambda(I)\neq 0\), then
      \[
      \mathrm{Ext}(I)\neq\varnothing;
      \]
\item if \(\Lambda(I)\neq 0\) and \(J\succeq I\) satisfies \(\Lambda(J)=0\),
      then \(J\in \mathrm{Ext}(I)\).
\end{enumerate}
\end{axiom}

\begin{corollary}[Nonterminal incidences have nonempty continuation class]
\label{cor:nonterminal-incidences-continue}\lean{CoherenceCalculus.nonterminal_incidences_continue}
If an incidence \(I\) is not terminal, then \(\Lambda(I)\neq 0\) and
\[
\mathrm{Ext}(I)\neq\varnothing.
\]
\end{corollary}

\begin{proof}
If \(I\) is not terminal, Definition~\ref{def:terminal-incidence} implies that
Constancy does not lift to \(I\). By
Axiom~\ref{ax:termination}, this is equivalent to \(\Lambda(I)\neq 0\), and
then Axiom~\ref{ax:nonretroactive-exactification}(i) yields
\(\mathrm{Ext}(I)\neq\varnothing\).
\end{proof}

\begin{corollary}[Zero defect cannot arise in place]
\label{cor:zero-defect-cannot-arise-in-place}\lean{CoherenceCalculus.zero_defect_cannot_arise_in_place}
If \(\Lambda(I)\neq 0\) and \(J\succeq I\) satisfies \(\Lambda(J)=0\), then
\[
J\succ I
\]
and \(J\) preserves the already-incident primitive relations of \(I\).
\end{corollary}

\begin{proof}
By Axiom~\ref{ax:nonretroactive-exactification}(ii), \(J\in\mathrm{Ext}(I)\). The
claim follows from Definition~\ref{def:exactifying-continuation-class}.
\end{proof}

\begin{remark}[Termination and continuation, not yet a chosen dynamics]
\label{rem:exactification-no-hidden-selector}\lean{CoherenceCalculus.TerminationData, CoherenceCalculus.NonretroactiveExactificationData, CoherenceCalculus.ExactifyingContinuationClass}
The termination law identifies the exact endpoint. The exactification law does
not yet choose a unique continuation
\[
I \mapsto J.
\]
It forces only a nonempty exactifying continuation class \(\mathrm{Ext}(I)\) whenever
\(I\) is nonterminal. Selection, minimality, and realized law choice enter
later through the admissibility and selected-cut layers. This prevents a hidden
selector from being smuggled into bedrock under the name of exactification.
Likewise, the present bedrock laws identify the zero-defect endpoint and the
need for continuation, not yet a transport law for defect burden along
extension. Any stronger monotonicity or comparison statement for unresolved
content must be added later as an explicit theorem, not inferred rhetorically
from the existence of \(\mathrm{Ext}(I)\).
\end{remark}

\begin{remark}[No physical time has been introduced]
\label{rem:exactification-not-yet-time}\lean{CoherenceCalculus.TerminationData, CoherenceCalculus.ExactifyingContinuationClass, CoherenceCalculus.NonretroactiveExactificationData}
The termination/exactification laws are bedrock closure and continuation laws,
not yet temporal ones. They say only that terminality is exact zero defect,
and that nonterminal incidence cannot close in place and therefore must
continue by strict extension. Ordered persistence, temporal interface, and
physical before/after are later realization layers built on top of these
bedrock laws.
\end{remark}

% ------------------------------------------------------------
